\section{Related work}
\label{sec:related}

\paragraph{Decoding auditory attention.}
Brain activity follows the speech a listener attends to \citep{ding2012,lalor2009}.
The oldest decoders are linear: they rebuild the speech envelope from the EEG and
compare the reconstruction with each talker \citep{osullivan2015}. They are
simple and well understood, but they need windows of tens of seconds to decide
reliably \citep{geirnaert2021}, too slow for a hearing device. Deep models
shorten the window in two ways. \emph{Match-mismatch} models learn directly
whether a window of EEG goes with a window of speech
\citep{accou2021,monesi2020,puffay2023}; our decoder belongs to this family.
\emph{Direction} decoders drop the audio and classify which side the listener
attends
\citep{vandecappelle2021,yan2024darnet,fan2025listennet,ni2024dbpnet,xu2024densenet,zhang2025swim}.
Both lines share a weakness: the reported accuracy often measures something other
than attention. \citet{rotaru2024} showed that direction decoding on several
public datasets relies largely on eye movements captured in the EEG rather than
on attention, and when within-trial leakage between training and test is removed,
recent direction decoders lose $17$ to $45$ points and several fall to chance
\citep{alavi2026neurotoken} (Table~\ref{tab:published}). We bring the concern of
\citet{rotaru2024} to the matching task and to several recorded streams at once,
and make the check part of training and model choice rather than an
after-the-fact test.

\paragraph{Datasets.}
The public AAD corpora record EEG and little else, and their designs deliberately
hold the listener still, which is good for isolating the brain response and
rules out the multimodal question: DTU uses earphones that simulate a room
\citep{fuglsang2018dtu}, KU~Leuven a quiet booth \citep{biesmans2017kul,das2019kuldata},
and a recent four-speaker dataset a chin rest \citep{yan2024fourspeaker}. A
stream that was never recorded, from a listener who was not free to move, cannot
compete with the EEG for credit. KU~Leuven has a further limitation we measure in
Appendix~\ref{app:external}: one of its two talkers is attended in
\R{extprior/kul/nrole} of \R{extprior/kul/ntrials} trials, so the audio alone
carries part of the answer. AV-GC-AAD \citep{geirnaert2024avgc} records gaze and
decouples it from attention on purpose --- the right tool for asking
\emph{whether} a decoder reads gaze instead of the brain, but its other streams
are designed to carry no information, so it cannot say how credit is shared among
streams that do. MAESTRO \citep{maestro2026} records EEG, eye tracking, head
motion and scene video from a listener who moves freely in a room of
loudspeakers, which is why we work on it; its own limitation, an audio confound
in the usual candidate construction, is the subject of
Section~\ref{sec:failures}.

\paragraph{Shortcut learning and multimodal balance.}
The problem is a case of shortcut learning
\citep{geirhos2020,torralba2011,gururangan2018}. The usual fixes rebalance the
data \citep{goyal2017vqa} or train against a shortcut-only branch
\citep{cadene2019rubi}; we add a scoring function in which the shortcut is
impossible. Work on balancing modalities
\citep{wang2020gradblend,huang2022competition,wu2022greedy,peng2022ogm} evens out
how fast each stream learns during training; we measure each stream's effect on
the decision at test time instead. The measurement is permutation importance
\citep{breiman2001,fisher2019} and the Shapley value
\citep{shapley1953,lundberg2017,covert2020} applied to whole streams; what is new
is using it to choose models and to report every result.
Appendix~\ref{app:related} gives a longer discussion.
